\section{Einleitung}
\label{sec:einleitung}

Aluminium ist einer der wichtigsten metallischen Rohstoffe und wird in vielen Industriezweigen genutzt wird.
Die Gewinnung und der Handle von Aluminium erfolgt weltweit weswegen Aluminium an der London Metal Exchange Börse gehandelt wird.
Der Preis von Aluminium ist sehr Volatil da er von vielen Faktoren wie Energiepreise im Herstellungsland, Transportkosten oder Einfuhrzölle abhängt.
Ziel dieser Arbeit ist es die Prozentuale veränderung des Aluminiumpreises für einen Tag und bis zu drei Wochen vorhersagen.
Um dies zu erreichen werden verschiedene Machine Learning Modelle trainiert. 
Die Datengrundlage dafür bieten historische Preise von Aluminium und weiteren Metallen bei denen Daten frei zugänglich waren, sowie der Gesamtwirtschaftliche Index S\&P 500.
Die trainierten Modelle werden anschließend in verschiedenen Metriken verglichen und ausgewertet.

\subsection{Komplexität der Aufgabenstellung}
Die Modellierung des Aluminiumpreises ist eine komplexe Problemstellung, da der Preis von ökonomischen, geopolitischen und technologischen Faktoren beeinflusst wird. Neben Angebot und Nachfrage wirken Indikatoren wie Inflationsraten, Wechselkursschwankungen des US-Dollars sowie energiepolitische Rahmenbedingungen auf den Preis ein. Zudem können unvorhergesehene geopolitische Ereignisse zu abrupten Abweichungen in der Zeitreihe führen. Diese Eigenschaften bedingen ein niedriges Signal-Rausch-Verhältnis. Das bedeutet, dass die relevanten Informationen für die Prognose in einem großen Rauschen von irrelevanten Daten verborgen sind. Daher ist es notwendig, robuste Modelle zu entwickeln, die in der Lage sind, die zugrundeliegenden Muster zu erkennen und gleichzeitig mit der Unsicherheit und Volatilität des Marktes umzugehen.

\subsection{Zielsetzung des Projekts}
Das Ziel dieses Projekts ist die Entwicklung und Evaluation eines datengetriebenen Systems zur kurz- bis mittelfristigen Prognose des Aluminiumpreises. Hierbei kommen Methoden der Zeitreihenanalyse und des maschinellen Lernens zum Einsatz. Der Fokus liegt auf dem Vergleich klassischer statistischer Baselines mit modernen Deep-Learning-Architekturen. Ein zentraler Aspekt ist die Berücksichtigung von Unsicherheiten in den Vorhersagen, um eine bessere Entscheidungsgrundlage für die Produktionsplanung zu schaffen. Der gewählte Ansatz integriert dabei weitere Variablen anderer Metallmärkte sowie einen gesamtwirtschaftlichen Index, um die Vorhersage zu optimieren.


